\specsection{附录的习题}

\begin{enumerate}
\def\labelenumi{\arabic{enumi})}
\tightlist
\item
  设 E 是特征 \(\neq 2\) 的交换域上的一个有限维向量空间。用 H 表示 E 上一个非退化二次型，用 S 表示 \(E \times E\) 上使得对 E 中所有 x 有 \(\mathrm{H}(x) = \mathrm{S}(x, x)\) 的对称双线性型。设 Q 是 E 上的一个二次型。
\end{enumerate}

\begin{enumerate}
\def\labelenumi{\alph{enumi})}
\item
  存在自同态 \(u\)（\(\mathbf{E}\) 的），它由关系式 \(\mathbf{Q}(x) = \mathrm{S}(u(x), x)\) 与 \(\mathrm{S}(u(x), y) = \mathrm{S}(x, u(y))\)（\(x\)、\(y\) 在 \(\mathbf{E}\) 中）刻画。置 \(\operatorname{Tr}(\mathbf{Q}/\mathbf{H}) = \operatorname{Tr}(u)\)。
\item
  推广 No.~1 的 Remark 3。
\item
  若 \((e_i)_{1\leqslant i\leqslant m}\) 是 E 的对 H 标准正交的基，则 \(\operatorname{Tr}(\mathbf{Q}/\mathbf{H}) = \sum_{i=1}^{m}\mathrm{Q}(e_i)\)。
\end{enumerate}

\begin{enumerate}
\def\labelenumi{\arabic{enumi})}
\setcounter{enumi}{1}
\tightlist
\item
  设 E 是一个实 Hilbert 空间，Q 是 E 上一个连续正二次型，H 是 E 上的二次型 \(x \mapsto \|x\|^2\)。
\end{enumerate}

\begin{enumerate}
\def\labelenumi{\alph{enumi})}
\item
  证明 \(\operatorname{Tr}(\mathbf{Q}/\mathbf{H}) = \sum_{i \in I} \mathbf{Q}(e_i)\) 对 E 的每个标准正交基 \((e_i)_{i \in I}\) 成立。（设 \((a_1, \ldots, a_p)\) 是 E 中一个有限标准正交族；对每个 \(\varepsilon > 0\)，存在 I 的有限子集 J 与元素 \(a_1', \ldots, a_p'\)，它们是 \(e_i\)（对 \(i \in J\)）的线性组合，且满足 \(\|a_j - a_j' \| < \varepsilon\)（对 \(1 \leqslant j \leqslant p\)）；于是 \(\sum_{j=1}^{p} \mathbf{Q}(a_j') \leqslant \sum_{i \in J} \mathbf{Q}(e_j) \leqslant \sum_{i \in I} \mathbf{Q}(e_i)\)。由此推出 \(\sum_{j=1}^{p} \mathbf{Q}(a_j) \leqslant \sum_{i \in I} \mathbf{Q}(e_i)\)。）
\item
  由 \(a\) ) 推出 Prop. 3 的一个新证明。
\item
  设 \(\mathbf{E}_0\) 是 \(\mathbf{E}\) 的一个稠密线性子空间，\(\mathbf{Q}_0\)（相应地，\(\mathbf{H}_0\)）是 \(\mathbf{Q}\)（相应地，\(\mathbf{H}\)）在 \(\mathbf{E}_0\) 上的限制。证明 \(\operatorname{Tr}(\mathbf{Q}/\mathbf{H}) = \operatorname{Tr}(\mathbf{Q}_0/\mathbf{H}_0)\)。（设 \((e_1, \ldots, e_n)\) 是 \(\mathbf{E}\) 中一个线性无关的序列，生成一个子空间 \(\mathbf{F}\)。用 \(\mathbf{Q}_{\mathbf{F}}\)（相应地，\(\mathbf{H}_{\mathbf{F}}\)）表示 \(\mathbf{Q}\)（相应地，\(\mathbf{H}\)）在 \(\mathbf{F}\) 上的限制。利用 No.~1 的 Remark 3，证明 \(\operatorname{Tr}(\mathbf{Q}_{\mathbf{F}}/\mathbf{H}_{\mathbf{F}})\) 是 \((e_1, \ldots, e_n)\) 的连续函数。）
\end{enumerate}

\begin{enumerate}
\def\labelenumi{\arabic{enumi})}
\setcounter{enumi}{2}
\item
  设 E 与 F 是两个实向量空间，u 是 E 到 F 中的一个线性映射；若 Q 与 H 是 F 上的正二次型，使得 \(\mathrm{H}(x)=0\) 蕴含 \(\mathrm{Q}(x)=0\)（对 \(x\in F\)），则 \(\mathrm{Tr}\left(\mathrm{Q}\circ\mathrm{u}/\mathrm{H}\circ\mathrm{u}\right)\leqslant\mathrm{Tr}\left(\mathrm{Q}/\mathrm{H}\right)\)。
\item
  设 I 是一个集合，\(l^2(\mathbf{I})\) 是由实数族 \(\mathbf{x} = (x_i)_{i \in \mathbf{I}}\)（满足 \(\sum_{i \in \mathbf{I}} x_i^2\) 有限）组成的向量空间。设 \((\lambda_i)_{i \in \mathbf{I}}\) 是一个可和的正数族。对每个 \(\mathbf{x}\)（在 \(l^2(\mathbf{I})\) 中），置 \(Q(\mathbf{x}) = \sum_{i \in \mathbf{I}} \lambda_i x_i^2\) 与 \(H(\mathbf{x}) = \sum_{i \in \mathbf{I}} x_i^2\)。证明 \(\operatorname{Tr}(Q / H) = \sum_{i \in \mathbf{I}} \lambda_i\)。（利用 Exer. 2 a)。
\item
  设 E 是一个实向量空间，\((\lambda_i)_{i \in I}\) 是一个可和的正数族，\((y_i)_{i \in I}\) 是 E 上线性形式组成的一个族。对每个 \(x \in E\)，置 \(\mathrm{H}(x) = \sum_{i \in I} \langle x, y_i \rangle^2\) 与 \(\mathrm{Q}(x) = \sum_{i \in I} \lambda_i \langle x, y_i \rangle^2\)；假设 \(\mathrm{H}(x)\)（对每个 \(x \in E\)）有限。证明 \(\mathrm{Q}(x)\)（对每个 \(x \in E\)）有限，Q 与 H 是 E 上的正二次型，且 \(\operatorname{Tr}(\mathrm{Q}/\mathrm{H}) \leqslant \sum_{i \in I} \lambda_i\)。（置 \(u(x) = (\langle x, y_i \rangle)_{i \in I}\)，并对线性映射 \(u\)（E 到 \(l_2(I)\) 中的）应用 Exer. 3。）
\item
  设 E 是一个实向量空间，Q、H 与 \(H_{0}\) 是 E 上的正二次型。假设 \(H \leqslant a \cdot H_{0}\)，其中 a 是一个正实数。证明 \(\operatorname{Tr}(Q/H_{0}) \leqslant a \cdot \operatorname{Tr}(Q/H)\)。（利用 No.~1 的 Remark 1 化为 E 有限维的情形，然后利用 E 的一个对 H 与 \(H_{0}\) 都正交的基的存在性，由 Prop. 1 得出结论。）
\item
  设 E 是一个实 Hilbert 空间。证明 E 上的 Sazonov 拓扑由半范数 \(Q^{1/2}\) 组成的集合定义，其中 Q 取遍 E 上核正二次型所成的集合。（利用 Exer. 6。）
\item
  设 E 与 F 是两个实 Hilbert 空间。在 \(E \otimes F\) 上存在一个 Hausdorff 准 Hilbert 空间结构，使得 \((x_{1} \otimes y_{1}|x_{2} \otimes y_{2}) = (x_{1}|x_{2}) \cdot (y_{1}|y_{2})\)（对 E 中的 \(x_{1}, x_{2}\) 与 F 中的 \(y_{1}, y_{2}\)）。用 \(E \otimes_{2} F\) 表示 \(E \otimes F\) 的 Hilbert 空间完备化。
\end{enumerate}

\begin{enumerate}
\def\labelenumi{\alph{enumi})}
\item
  若 \(\mathbf{E}'\)（相应地，\(\mathbf{F}'\)）是 \(\mathbf{E}\)（相应地，\(\mathbf{F}\)）的一个闭线性子空间，证明 \(\mathbf{E}' \otimes_2 \mathbf{F}'\) 可以等同于 \(\mathbf{E} \otimes_2 \mathbf{F}\) 中由元素 \(x \otimes y\)（满足 \(x \in \mathbf{E}'\) 与 \(y \in \mathbf{F}'\)）生成的闭线性子空间。
\item
  假设 E（相应地，F）是闭线性子空间族 \((\mathrm{E}_{\alpha})_{\alpha\in\mathbf{A}}\)（相应地，\((\mathrm{F}_{\beta})_{\beta\in\mathbf{B}}\)）的 Hilbert 和（TVS, V, §2, No.~2, Def. 2）。证明 \(E\otimes_{2}F\) 是闭线性子空间族 \((\mathrm{E}_{\alpha}\otimes_{2}\mathrm{F}_{\beta})_{(\alpha,\beta)\in\mathbf{A}\times\mathbf{B}}\) 的 Hilbert 和。
\item
  若 \((e_i)_{i\in \mathbf{I}}\)（相应地，\((f_j)_{j\in \mathbf{J}})\)）是 E（相应地，F）的标准正交基，则族 \((e_i\otimes f_j)_{(i,j)\in \mathbf{I}\times \mathbf{J}}\) 是 \(\mathbf{E}\otimes_2\mathbf{F}\) 的标准正交基。
\item
  设 \(\mathbf{G}\) 是一个实 Hilbert 空间。定义从 \(\mathbf{E} \otimes_{2} \mathbf{F}\) 到 \(\mathbf{F} \otimes_{2} \mathbf{E}\) 上以及从 \((\mathbf{E} \otimes_{2} \mathbf{F}) \otimes_{2} \mathbf{G}\) 到 \(\mathbf{E} \otimes_{2} (\mathbf{F} \otimes_{2} \mathbf{G})\) 上的典范同构。
\item
  设 \(E_{1}\) 与 \(F_{1}\) 是两个实 Hilbert 空间，\(u:E\to E_{1}, v:F\to F_{1}\) 是连续线性映射。证明 \(u\otimes v\) 可以由连续性延拓为连续线性映射 \(u\otimes_{2}v\)（\(E\otimes_{2}F\) 到 \(E_{1}\otimes_{2}F_{1}\) 中的），且 \(\|u\otimes_{2}v\|=\|u\|\cdot\|v\|\)。
\end{enumerate}

\begin{enumerate}
\def\labelenumi{\arabic{enumi})}
\setcounter{enumi}{8}
\tightlist
\item
  设 E 与 F 是两个实 Hilbert 空间。
\end{enumerate}

\begin{enumerate}
\def\labelenumi{\alph{enumi})}
\item
  证明存在一个连续线性映射 \(\varphi\)（\(\mathbf{E} \otimes_2 \mathbf{F}\) 到 \(\mathcal{L}(\mathbf{E}; \mathbf{F})\) 中的），它由下列条件刻画：\(\varphi(x \otimes y)(x') = (x|x') \cdot y\)（对 \(x, x'\)（\(\mathbf{E}\) 中的）与 \(y\)（\(\mathbf{F}\) 中的））。证明 \(\varphi\) 的范数为 1，且 \(\varphi(\mathbf{E} \otimes \mathbf{F})\) 是 \(\mathbf{E}\) 到 \(\mathbf{F}\) 中的有限秩连续线性映射所成的集合。
\item
  证明 \(\varphi\) 是 \(E \otimes_2 F\) 到 E 到 F 中的 Hilbert--Schmidt 线性映射所成的集合 HS (E, F) 上的一个双射（利用 Exer. 8 c) 与 Prop. 3）。给 HS (E, F) 赋以从 \(E \otimes_2 F\) 的 Hilbert 空间结构借助双射 \(\varphi\) 导出的结构；相应的范数记作 \(\|u\|_2\)。设 \(u \in \mathrm{HS}(E, F)\)；定义 E 上的正二次型 \(Q_u\) 与 H 如下：\(Q_u(x) = \|u(x)\|^2\) 与 \(H(x) = \|x\|^2\)。证明 \(\|u\|_2^2 = \operatorname{Tr}(Q_u / H)\)。
\item
  设 \(\mathbf{E}_1\) 与 \(\mathbf{F}_1\) 是两个实 Hilbert 空间，\(u: \mathbf{E}_1 \to \mathbf{E}\)、\(v: \mathbf{F} \to \mathbf{F}_1\) 是连续线性映射。设 \(\varphi_1\) 是 \(\mathbf{E}_1 \otimes_2 \mathbf{F}_1\) 到 \(\mathrm{HS}(\mathbf{E}_1, \mathbf{F}_1)\) 上的按类似于 \(\varphi\) 的方式定义的同构。证明 \(v \circ \varphi(t) \circ u = \varphi_1((u^* \otimes v)(t))\) 对所有 \(t \in \mathbf{E} \otimes_2 \mathbf{F}\) 成立。由此推出，对每个 Hilbert-Schmidt 映射 \(w\)（\(\mathbf{E}\) 到 \(\mathbf{F}\) 中的），线性映射 \(v \circ w \circ u\)（\(\mathbf{E}_1\) 到 \(\mathbf{F}_1\) 中的）是 Hilbert-Schmidt 的，且 \(\|v \circ w \circ u\|_2 \leqslant \|v\| \cdot \|w\|_2 \cdot \|u\|\)。
\end{enumerate}

